\documentclass{beamer}
%\documentclass{article}
%\usepackage{beamerarticle}
\mode<article>{\usepackage{fullpage}}

\usepackage{polyglossia}
\setmainlanguage{italian}
\setotherlanguage{english}

\PolyglossiaSetup{italian}{indentfirst=false}
\mode<presentation>{\usetheme{moloch}}
\mode<presentation>{\useinnertheme{moloch}}
\mode<presentation>{\usefonttheme{moloch}}
\mode<presentation>{\setbeamertemplate{itemize items}[triangle]}

%\usecolortheme{seahorse} 
\title{Worspace 365 Canva}
\author{\texorpdfstring{John Doe\newline\url{info@iisperugia.net}}{John Doe}}
%\author{}
\begin{document}
\begin{frame}%[plain]
    \maketitle
 \end{frame}
\begin{frame}{Sommario}
	\tableofcontents
\end{frame}
\section{Introduzione}
\begin{frame}{Quattro chiacchiere}
	La nostra scuola ha 3 ambienti 
	\begin{itemize}[(I)]
		\item<1-> Google Workspace for Education Fundamentals
		\item<2-> Microsoft 365
		\item<3-> Canva
	\end{itemize}
\end{frame}
\section{Google Workspace}
\begin{frame}{Google Workspace}
 Google Workspace ha tre compiti
	\begin{itemize}[(I)]
		\item<1-> Gestire la posta
		\item<2-> Gsiute 
		\item<3-> SSO
		\begin{itemize}
			\item<4->Microsof 365
			\item<5->Canva
		\end{itemize}
	\end{itemize}
\end{frame}
\section{SSO}
\begin{frame}{SSO}
	\begin{alertblock}{SSO}
	Single Sign-On (SSO) è uno strumento di autenticazione che consente agli utenti di accedere in modo sicuro a più applicazioni e servizi utilizzando un unico set di credenziali, eliminando la necessità di ricordare password diverse per ciascun servizio\footnote{Informazione ricavata da https://www.okta.com/blog/2021/02/single-sign-on-sso/ consultato il 27/07/2024}.
	\end{alertblock}
\end{frame}
\begin{frame}{Unico utente e unica password}
	\begin{alertblock}{Unico utente e unica password}
	nome.cognome.doc@iisperugia.net
	
	password 
	
Danno accesso a tutte le piattaforme della scuola
	\end{alertblock}
\end{frame}
\section{Sotto il cofano}
\subsection{Perché dobbiamo cancellare gli utenti?}
\begin{frame}{Perché dobbiamo cancellare gli utenti?}
	\begin{itemize}
		\item<1-> Garante della Privacy
		\item<2-> Obblighi contrattuali
		\item<3-> Numero di licenze disponibili
	\end{itemize}
\end{frame}
\begin{frame}{Garante della Privacy}
\begin{exampleblock}{Garante della Privacy}
Abbiamo in deposito documenti personali e quindi, alla cessazione, tali documenti devono essere distrutti.
\end{exampleblock}
\end{frame}
\begin{frame}{Obblighi contrattuali}
	\begin{block}{Obblighi contrattuali}
	La scuola può fornire questi servizi solo a studenti e personale scolastico che appartengono a questa scuole.
	\end{block}
\end{frame}
\begin{frame}{Numero di licenze disponibili}
	\begin{alertblock}{Numero di licenze disponibili}
		Il problema è Canva che ci ha fornito solo di un numero limitato di licenze.
	\end{alertblock}
\end{frame}
\subsection{Proprietario}
\begin{frame}[standout]{Proprietario}
	\begin{block}{Proprietario}
	Qualunque utente che crei qualcosa in WorkSpace è un proprietario
	\end{block}
\end{frame}
\begin{frame}{Proprietario e gli altri}
	\begin{block}{Proprietario e gli altri}
	Un Proprietario concede agli altri le sue cose impostando vari livelli di accesso.
	\end{block}
	\begin{alertblock}{Accesso}
		Con i livelli di accesso non viene trasferita la proprietà  
	\end{alertblock}
	\begin{exampleblock}{Esistenza}
		Il documento, drive, etc. esiste finché  l'utente esiste.
	\end{exampleblock}
\end{frame}
\subsection{Drive personali Vs Drive condivisi}
\begin{frame}{Drive personali Vs Drive condivisi}
	\begin{alertblock}{Drive personali}
	 Condivido con gli altri ma mantengo la proprietà
	\end{alertblock}
	\begin{block}{Drive Condivisi}
		Condivido con gli altri ma perdo la proprietà
	\end{block}
\end{frame}
\section{Navigazione Incognita}
\begin{frame}{Navigazione Incognita}
\centering
\includegraphics[width=0.50\textwidth]{navinc}
\end{frame}
\begin{frame}[standout]
\centering
\Huge\bfseries
\textcolor{orange}{The End}
\end{frame}
\end{document}
